% TOP_PROBLEM: {"id": 30006899, "title": "Pugh–Shub Stable Ergodicity (C^r Density in Partial Hyperbolicity)", "category_id": 11, "category": {"id": 11, "name": "dynamical_systems", "display_name": "Dynamical Systems", "description": "Problems about long-term behavior of deterministic systems, Hamiltonian dynamics, and periodic orbits.", "slug": "dynamical-systems", "order_index": 11, "created_at": "2026-07-31T15:26:25.671Z"}, "status": "open"}
% FIRST_POSTED: 2026-09-27
% !TeX program = lualatex
% ATTEMPT_STATUS: unresolved
% Model: GPT-6 Astra Ultra; continuation dated 27 September 2026.
\documentclass[11pt]{article}
\usepackage[margin=25mm]{geometry}
\usepackage{fontspec}
\setmainfont{Latin Modern Roman}
\usepackage{amsmath,amssymb,mathtools}
\usepackage{unicode-math}
\setmathfont{Latin Modern Math}
\usepackage{xurl,hyperref}
\hypersetup{colorlinks=true,urlcolor=blue}
\setcounter{secnumdepth}{0}
\setlength{\parindent}{0pt}
\setlength{\parskip}{5pt}
\setlength{\emergencystretch}{3em}
\begin{document}
\section*{\texorpdfstring{Pugh–Shub Stable Ergodicity (C\textasciicircum{}r Density in Partial Hyperbolicity)}{Pugh–Shub Stable Ergodicity (C\textasciicircum{}r Density in Partial Hyperbolicity)}}\label{30006899}
\subsection{Definitions and mathematical statement}
Let $(M,m)$ be a compact connected boundaryless smooth manifold with smooth probability volume. Partial hyperbolicity is a continuous invariant splitting $TM=E^u\oplus E^c\oplus E^s$ with nonzero $E^u,E^s$, possible zero center, and a common iterate $N$ such that unstable unit vectors expand by more than two, stable ones contract below one-half, and earlier nonzero bundles dominate later ones by a factor greater than two at the same base point. A map $g$ is stably ergodic here if a $C^1$ neighborhood of it has every $C^2$ volume-preserving member ergodic for $m$, meaning every invariant measurable set has measure zero or one. For every $r\in[2,\infty]$, the assertion is density of these $g$ in $\operatorname{PH}_m^r(M)$ in the relative $C^r$ topology. Nonintegral finite $r$ uses H\"older regularity. The density topology and the stability topology are intentionally different.
\par\smallskip\textit{Formulation and concept references:} \href{https://arxiv.org/pdf/1709.04983}{[S1]}.

\subsection{Short English statement}
Can every volume-preserving partially hyperbolic map be approximated in its full $C^r$ topology by one whose nearby $C^2$ conservative maps remain ergodic?
\subsection{Research attempt}
\paragraph{Two different topologies occur in the question.}
The inspected source [S1, Theorem A and Theorem A'] proves
$C^1$ density of stable ergodicity among the stated partially
hyperbolic conservative maps. This is substantial progress, but
the selected assertion asks for approximation in the relative
$C^r$ topology, for every $r\geq2$, while retaining $C^1$ stability
against $C^2$ conservative perturbations. We do not upgrade the
source's density topology by changing the regularity label of the
approximating maps.

The partial results below give explicit hyperbolic examples,
prove ergodicity but not stable ergodicity in a natural center
family, and identify the quantitative perturbation and accessibility
issues left by that family. The source's deep perturbation and
foliation results are kept distinct from the elementary Fourier
and conjugacy calculations supplied here.

\paragraph{An explicit ergodic uniformly hyperbolic map.}
On $\mathbb T^2=\mathbb R^2/\mathbb Z^2$, consider
\[
 A=\begin{pmatrix}2&1\\1&1\end{pmatrix}.
\]
It has determinant one and eigenvalues
$\lambda=(3+\sqrt5)/2>1$ and $\lambda^{-1}<1$.
The constant eigendirections provide a stable and unstable
splitting, and some common iterate meets the numerical inequalities
in the definition. The map preserves normalized area because
its determinant is one.

Here is a direct proof of ergodicity. If $f\in L^2(\mathbb T^2)$
is invariant, its Fourier coefficients are constant along every
orbit of $A^{\mathsf T}$ on $\mathbb Z^2$. A nonzero integer
vector cannot have a finite orbit: that would give
$((A^{\mathsf T})^k-I)v=0$ for some $k>0$, whereas neither
$\lambda^k$ nor $\lambda^{-k}$ is one. Every nonzero orbit is
therefore infinite. Square summability of Fourier coefficients
forces the constant coefficient on each such orbit to be zero.
Thus every invariant $L^2$ function is constant. Applying this
to invariant-set indicators proves ergodicity.

This argument proves the single map's ergodicity without an
unstated appeal to topological transitivity. Stable ergodicity
requires additional control of nonlinear perturbations.

\paragraph{The Hopf criterion and what its use imports.}
For a $C^2$ conservative Anosov diffeomorphism, the stable and
unstable foliations are absolutely continuous and have local
product structure. We use these standard hyperbolic-dynamics
facts as external input, not as consequences of the Fourier proof.
They are the kind of foliation input explicitly discussed in
[S1]. Given them, the ergodicity argument can be described precisely.

Take a continuous observable $\phi$. Birkhoff's theorem gives
forward and backward averages on a full-measure set, and these
averages agree almost everywhere, as both equal conditional
expectation onto the invariant sigma-algebra. Points on the same
stable leaf have the same forward average whenever defined: their
forward distances tend to zero, so uniform continuity of $\phi$
makes the difference of the Cesaro averages tend to zero. Points
on an unstable leaf have the same backward average by the same
reason in reverse time.

Absolute continuity allows these full-measure statements to pass
to almost every leaf plaque. In a local product box, a function
that is essentially constant on both transverse plaque families
is constant almost everywhere: in product coordinates it is a
function of the first coordinate alone and also of the second
alone, and Fubini forces a common constant. The local product
measure is equivalent to the smooth volume measure by the stated
foliation input. Overlapping boxes have the same constant, and
connectedness of the manifold gives one global constant.

Applying this to a countable dense family of continuous observables
shows that their invariant conditional expectations are constants.
Density then gives the same for all $L^2$ observables, proving
ergodicity. The step invoking absolute continuity is important:
topological accessibility of individual points alone would not
justify the passage through almost-everywhere Birkhoff sets.

The Anosov property is $C^1$ open. For the displayed linear map,
this can be seen from invariant cones around the two eigenlines:
in an adapted norm the expansion/contraction factors are
$\lambda,\lambda^{-1}$, and sufficiently small derivative
perturbations preserve slightly wider cones with strict factors.
Iteration of the cone graph maps produces invariant stable and
unstable subbundles. Thus every sufficiently $C^1$-close $C^2$
conservative map is again Anosov, and the preceding Hopf argument
proves its ergodicity. This verifies stable ergodicity of this
example in the exact stability convention of the question.

More generally the same standard Anosov inputs settle the
zero-center subcase. The difficulty in the catalogue question
lies in permitting a genuine center bundle, where strong stable
and unstable plaques no longer form a full-dimensional local
product chart.

\paragraph{A center family that is ergodic but never stably ergodic.}
On $\mathbb T^2\times\mathbb T$, let
\[
 F_\alpha(x,\theta)=(Ax,\theta+\alpha).
 \tag{388.1}
\]
Its derivative has the splitting into the expanding and contracting
eigenlines of $A$ and the vertical center line. The center multiplier
is one. A common large iterate has unstable expansion greater
than two, stable contraction less than one half, and the required
domination over and under the center. Thus it belongs to the
selected partially hyperbolic class for every $\alpha$ and is
smooth and volume preserving.

If $\alpha$ is irrational, $F_\alpha$ is ergodic. For an invariant
$L^2$ function, its coefficients indexed by $(k,\ell)\in
\mathbb Z^2\times\mathbb Z$ have constant moduli along the
$A^{\mathsf T}$ orbit of $k$, with a unit complex phase from
$e^{2\pi i\ell\alpha}$. For $k\ne0$ the orbit is infinite,
so square summability again forces vanishing. At $k=0$, invariance
requires the coefficient to equal its multiple by
$e^{2\pi i\ell\alpha}$, which forces it to vanish for
$\ell\ne0$. Only the constant coefficient remains.

If $\alpha=a/q$ is rational, the nonconstant function
$e^{2\pi iq\theta}$ is invariant, so the map is not ergodic.
Choose rationals $\alpha_j\to\alpha$. Then
$F_{\alpha_j}\to F_\alpha$ in every $C^r$ topology, including
$C^\infty$, because the difference is only a constant translation
in the center coordinate. Therefore no map in the family (388.1)
is stably ergodic, even when it is ergodic.

This is not a counterexample to density in the whole partially
hyperbolic space: perturbations outside this product family may
be stably ergodic. It proves that proving ergodicity on a dense
set of rotation parameters is much weaker than the desired
stable-ergodicity assertion.

\paragraph{The obstruction is visible in strong accessibility.}
For (388.1), stable and unstable leaves remain in a fixed circle
coordinate. A path made of finitely many strong stable or strong
unstable pieces therefore cannot change $\theta$. In the base,
the two eigendirections span the tangent plane; concatenating
pieces connects points in the same horizontal torus. Thus the
accessibility classes are precisely
$\mathbb T^2\times\{\theta\}$.

The simple Hopf argument above can at most show that an invariant
average is constant along such a class. It has no mechanism to
identify averages on different center fibers. For irrational
rotations the Fourier argument supplies ergodicity by a separate
mechanism, but rational approximation shows that mechanism is
not stable. Hence this family provides a concrete test for a
proof that claims to deduce stable ergodicity from expansion and
contraction alone.

\paragraph{A nonlinear-looking perturbation can still be a product in disguise.}
Let $h:\mathbb T^2\to\mathbb R$ be smooth and set
$G_h(x,\theta)=(x,\theta+h(x))$. It is a smooth volume-preserving
diffeomorphism. Direct calculation gives
\[
 G_hF_\alpha G_h^{-1}(x,\theta)
      =(Ax,\theta+\alpha+h(Ax)-h(x)).
 \tag{388.2}
\]
Thus a nonconstant fiber cocycle of the form
$\phi(x)=\alpha+h(Ax)-h(x)$ does not create genuinely new
ergodic behavior. It is conjugate to the product. Its stable and
unstable geometry may look tilted in the original coordinates,
but the invariant accessibility foliation is just the image of
the horizontal tori under $G_h$.

Such a skew product is still not stably ergodic: replacing
$\alpha$ by a nearby rational while keeping $h$ fixed gives
arbitrarily smooth-close nonergodic conjugate maps. For small $h$
it is an explicit perturbation within the partially hyperbolic
neighborhood of the product. More generally partial hyperbolicity
is invariant under this smooth conjugacy after taking an appropriate
common iterate and equivalent metric.

For any periodic base point $A^q x=x$, a cocycle of this form
satisfies the telescoping identity
\[
 \sum_{j=0}^{q-1}(\phi(A^jx)-\alpha)=0.
 \tag{388.3}
\]
Producing a perturbation that violates such an identity for a
chosen periodic orbit rules out this particular coboundary
description. It does not by itself prove accessibility or stable
ergodicity, but it is a useful necessary stress test for a
proposed construction. A proof cannot claim to break the product
obstruction while using only perturbations of the form (388.2).

\paragraph{Why C1-small perturbations do not solve the selected approximation.}
The topology issue is quantitative. On a torus coordinate strip,
choose a smooth bump $\psi$ supported in a short interval and,
for a small scale $\varepsilon$, form the volume-preserving shear
\[
 S_\varepsilon(x,y)=(x+\varepsilon^r\psi(y/\varepsilon),y),
 \tag{388.4}
\]
with the bump periodically extended from its small support. Its
Jacobian determinant is identically one. For an integer $r\geq2$,
its displacement tends to zero in $C^1$ at rate at most
$\varepsilon^{r-1}$, while its $r$-th derivative contains
$\psi^{(r)}(y/\varepsilon)$ and need not tend to zero at all.
At noninteger $r=m+\alpha$, $0<\alpha<1$, the amplitude
$\varepsilon^{m+\alpha}$ likewise keeps the $\alpha$-Holder
seminorm of the $m$-th derivative at constant scale while the
$C^1$ displacement tends to zero.

These are explicit conservative perturbations showing that a
$C^1$ error bound does not control the prescribed higher topology.
Making a perturbation smooth does not change this fact. The source's
theorem supplies smoothness of the approximants in an appropriate
class but its closeness conclusion is still only in $C^1$.
In the $C^\infty$ topology, each neighborhood controls a finite
collection of derivative seminorms; an argument must meet whichever
finite collection is prescribed, not just the first one.

One could try to spread a localized perturbation over a larger
region to reduce its high derivatives. That may alter the orbit
segments and invariant geometry it was meant to control. The
construction would need simultaneous estimates for the desired
ergodic mechanism and every derivative specified by the neighborhood.
No such estimate is supplied by the cone calculation or the
product-cocycle algebra above.

\paragraph{Diagnostics and outcome.}
The finite script checks the cat-map determinant and trace,
numerically checks the telescoping identity on sample periodic
orbits, and records the derived derivative-scaling exponents of
the conservative shear. The conjugacy is proved directly above.
These tests audit the concrete examples;
they do not certify a perturbation theorem for arbitrary center
dimension.

The proved elementary results distinguish ergodicity, stable
ergodicity and the two approximation topologies. With the stated
standard Anosov foliation input, the zero-center subcase is settled.
For a general partially hyperbolic map, the first missing step is
a perturbation in an arbitrary prescribed $C^r$ neighborhood that
creates a robust ergodic mechanism across the center directions.
The $C^1$ density theorem is not that step. The selected full
$C^r$ density assertion remains unresolved by this attempt.
\subsection{Sources}
\begin{itemize}
\item[S1]Pugh–Shub Stable Ergodicity (C\textasciicircum{}r Density in Partial Hyperbolicity); exact editorial selection using the primary formulations and hypotheses recorded in the source receipts.
\url{https://arxiv.org/pdf/1709.04983}.
\textit{Formulation source.}
\end{itemize}
\end{document}
